Many activities in modern daily life involve prolonged sitting, such as desk work, study, and commuting. Prolonged sitting with little lower-trunk movement is relevant because sitting pattern, not only total sitting duration, can affect how static a seated posture becomes. This thesis addresses a technical feasibility question: whether a low-cost pelvis-worn wearable prototype can be designed to monitor pelvic-region movement, collect labelled field data, run a compact local detector, and provide an on-device haptic cue during prolonged static sitting.

The prototype combines a Seeed Studio XIAO nRF54L15 Sense worn near the sacrum, the board's embedded LSM6DS3TR-C inertial measurement unit, Bluetooth Low Energy data streaming, an Android application for labelled field recording, a PC receiver for researcher-side inspection, a DRV2605L haptic driver, and a coin vibration motor. Four linked outcomes are evaluated. First, the thesis reports the implemented prototype and data-acquisition pipeline. Second, it evaluates detection of instructed pelvic rotation during train-sitting segments. In a person-independent split using participants P1--P4 for training and participant P5 for final testing, a \(k=7\) k-nearest-neighbours classifier achieved an F1-score of 0.957 for pelvic rotation, detecting 155 of 165 pelvic-rotation windows and producing 4 false positives among 3{,}497 static-sitting windows. Third, it compares the custom XIAO-based wearable with a research-grade Shimmer IMU at trend level. In a co-located comparison, Pearson correlations across acceleration and gyroscope axes ranged from 0.816 to 0.923 with estimated lag between \(-0.020\) and 0~s. In the valid field-worn paired recordings from two pilot participants, maximum cross-correlations reached 0.915--0.981 after lag correction. Finally, a compact deployable KNN with 12 features and 128 KMeans prototypes achieved an F1-score of 0.959 on the same held-out participant. On the XIAO nRF54L15 Sense, the compact feature-extraction and KNN cycle required 2.73~ms per 2~s window, using 56.3~KB of flash and 9.3~KB of RAM. The firmware implements a 10~min inactivity timer and a 1~s vibration reminder.

The findings show that the proposed low-cost wearable can connect pelvis-region IMU sensing, labelled mobile data collection, compact embedded classification, and local haptic cue delivery in one research prototype. The evidence supports trend-level sensing and instructed-movement detection in a preliminary pilot setting. It does not establish clinical pelvic kinematics, daily-life behavioural effectiveness, battery lifetime, or long-term wearing robustness.

\vspace{1em}
\underline{Keywords:} embedded machine learning, haptic feedback, inertial measurement unit, pelvic rotation, prolonged sitting, wearable sensing.
